% ═══════════════════════════════════════════════════════════════
% LH-01b: Newton 3 + Lageenergie – LEHRERHINWEISE
% Physik Klasse 10 MR
% ═══════════════════════════════════════════════════════════════

\documentclass[11pt, a4paper]{article}

\usepackage[utf8]{inputenc}
\usepackage[T1]{fontenc}
\usepackage[ngerman]{babel}

\usepackage{helvet}
\renewcommand{\familydefault}{\sfdefault}

\usepackage{microtype}
\usepackage[margin=2cm]{geometry}
\usepackage{enumitem}
\usepackage{tabularx}
\usepackage{booktabs}
\usepackage{array}
\usepackage{longtable}

\usepackage{amsmath, amssymb}
\usepackage{siunitx}
\sisetup{locale = DE, per-mode = symbol}

\usepackage{fancyhdr}
\usepackage{lastpage}

\usepackage{xcolor}
\usepackage{tcolorbox}
\tcbuselibrary{skins}

% ═══════════════════════════════════════════════════════════════
% FARBEN
% ═══════════════════════════════════════════════════════════════

\definecolor{physik}{HTML}{2c5aa0}
\definecolor{hellgrau}{HTML}{F5F5F5}
\definecolor{mittelgrau}{HTML}{888888}
\definecolor{loesung}{HTML}{c62828}
\definecolor{hinweis}{HTML}{b35c00}

\colorlet{fachfarbe}{physik}

% ═══════════════════════════════════════════════════════════════
% VARIABLEN
% ═══════════════════════════════════════════════════════════════

\newcommand{\fach}{Physik}
\newcommand{\thema}{Newton 3 + Lageenergie}
\newcommand{\klasse}{10 MR}
\newcommand{\stession}{01b}
\newcommand{\dauer}{90 Min (DS)}

% ═══════════════════════════════════════════════════════════════
% KOPF- UND FUSSZEILE
% ═══════════════════════════════════════════════════════════════

\pagestyle{fancy}
\fancyhf{}
\renewcommand{\headrulewidth}{0.4pt}
\renewcommand{\footrulewidth}{0pt}
\lhead{\small \fach{} Klasse \klasse{} | \thema{} | \textbf{Lehrerhinweise}}
\rhead{\small \dauer}
\cfoot{\small\textcolor{mittelgrau}{Seite \thepage{} von \pageref{LastPage}}}

% ═══════════════════════════════════════════════════════════════
% BOXEN
% ═══════════════════════════════════════════════════════════════

\newtcolorbox{infobox}[1][]{
    colback=hellgrau,
    colframe=fachfarbe,
    coltitle=white,
    colbacktitle=fachfarbe,
    boxrule=1pt,
    arc=0pt,
    left=8pt, right=8pt, top=6pt, bottom=6pt,
    before skip=8pt, after skip=8pt,
    fonttitle=\bfseries,
    title=#1
}

\newtcolorbox{warnbox}[1][]{
    colback=white,
    colframe=hinweis,
    coltitle=white,
    colbacktitle=hinweis,
    boxrule=1pt,
    arc=0pt,
    left=8pt, right=8pt, top=6pt, bottom=6pt,
    before skip=8pt, after skip=8pt,
    fonttitle=\bfseries,
    title=#1
}

% ═══════════════════════════════════════════════════════════════
% BEFEHLE
% ═══════════════════════════════════════════════════════════════

\newcommand{\lsg}[1]{\textcolor{loesung}{\textbf{#1}}}

\setlength{\parindent}{0pt}
\setlength{\parskip}{6pt}

% ═══════════════════════════════════════════════════════════════
% DOKUMENT
% ═══════════════════════════════════════════════════════════════

\begin{document}

% ═══════════════════════════════════════════════════════════════
% HEADER
% ═══════════════════════════════════════════════════════════════

\begin{center}
    {\Large\bfseries\textcolor{fachfarbe}{Lehrerhinweise: \thema}}\\[0.3em]
    {\small \fach{} Klasse \klasse{} | Stunde \stession{} | \dauer}
\end{center}
\vspace{-0.3em}
\hrule
\vspace{0.8em}

% ═══════════════════════════════════════════════════════════════
% ÜBERSICHT
% ═══════════════════════════════════════════════════════════════

\section*{Übersicht}

\begin{tabularx}{\textwidth}{@{}lX@{}}
\textbf{Rahmenplan:} & RP RS/GES 2021, Kl. 10: Wechselwirkung + Energiegewinnung \\
\textbf{Vorwissen:} & Newton 1+2 (01a), Hubarbeit (Kl. 7) \\
\textbf{Materialien:} & LP-01b (digital), AB-01b (Print), Beamer \\
\textbf{BE gesamt:} & 20 BE (4 Aufgaben) \\
\end{tabularx}

% ═══════════════════════════════════════════════════════════════
% STUNDENVERLAUF
% ═══════════════════════════════════════════════════════════════

\section*{Stundenverlauf}

\subsection*{Teil 1: Newton 3 – Wechselwirkungsgesetz (0–45')}

\begin{tabularx}{\textwidth}{@{}c|l|X|c@{}}
\toprule
\textbf{Zeit} & \textbf{Phase} & \textbf{Inhalt} & \textbf{SF} \\
\midrule
0–5' & Einstieg & Bild: Zwei Skater stoßen sich ab – wer bewegt sich? & PL \\
5–10' & Problem & „Wenn ich gegen die Wand drücke, drückt die zurück?" & PL \\
10–20' & Erarbeitung & LP Schritt 1–4: Wechselwirkungsgesetz einführen & EA \\
20–30' & Übung & AB Aufgabe 1a–c: Kräftepaare zeichnen & EA \\
30–40' & Simulation & SIM Tauziehen: Gleiche F, verschiedene a & EA \\
40–45' & Sicherung & AB Kasten 1 ausfüllen, Partner-Check & PA \\
\bottomrule
\end{tabularx}

\subsection*{Teil 2: Lageenergie (45–90')}

\begin{tabularx}{\textwidth}{@{}c|l|X|c@{}}
\toprule
\textbf{Zeit} & \textbf{Phase} & \textbf{Inhalt} & \textbf{SF} \\
\midrule
45–50' & Überleitung & „Woher kommt Energie beim Fallen?" – Hubarbeit & PL \\
50–60' & Erarbeitung & LP Schritt 8–11: $E_\text{pot} = m \cdot g \cdot h$ & EA \\
60–70' & Übung & AB Aufgabe 2a–c: Lageenergie berechnen & EA \\
70–80' & Simulation & SIM Lageenergie: Slider für m, h & EA \\
80–85' & Vertiefung & AB Aufgabe 3+4: Umstellen, Staudamm & EA \\
85–90' & Sicherung & AB Kasten 2+3 prüfen, Zusammenfassung & PL \\
\bottomrule
\end{tabularx}

\vspace{0.5em}
\textit{SF: PL = Plenum, EA = Einzelarbeit, PA = Partnerarbeit}

% ═══════════════════════════════════════════════════════════════
% DIFFERENZIERUNG
% ═══════════════════════════════════════════════════════════════

\section*{Differenzierung}

\begin{tabularx}{\textwidth}{@{}c|l|X@{}}
\toprule
\textbf{Niveau} & \textbf{Aufgaben} & \textbf{Hinweise} \\
\midrule
$\star$ Basis & 1a, 2a & Nur Einsetzen, keine Umstellung; bei Bedarf Werte vorgeben \\
$\star\star$ Standard & 1b, 2b, 2c, 3 & Formel-Dreieck nutzen; Zwischenschritte zeigen \\
$\star\star\star$ Erweitert & 1c, 4 & Transfer auf neue Situationen; Begründung verlangen \\
\bottomrule
\end{tabularx}

\begin{warnbox}[Zeitpuffer]
Bei Zeitmangel: Aufgabe 4 (Staudamm) als Hausaufgabe geben.\\
Bei Zeitüberschuss: Zusatzfrage „Warum beschleunigt der leichtere Skater schneller?"
\end{warnbox}

% ═══════════════════════════════════════════════════════════════
% HÄUFIGE FEHLER
% ═══════════════════════════════════════════════════════════════

\section*{Häufige Schülerfehler}

\begin{tabularx}{\textwidth}{@{}l|X|X@{}}
\toprule
\textbf{Fehler} & \textbf{Diagnose} & \textbf{Reaktion} \\
\midrule
„Die Kräfte heben sich auf" & Vergisst: Kräfte an VERSCHIEDENEN Körpern & „An welchem Körper wirkt $F_1$? An welchem $F_2$?" \\
\midrule
„Der Stärkere gewinnt" & Verwechselt Kraft mit Beschleunigung & SIM zeigen: Gleiche F, aber $a = F/m$ unterschiedlich \\
\midrule
$g = 9{,}81$ verwenden & Falsch gelernter Wert & „Wir rechnen IMMER mit $g = 10$ m/s² (kopfrechenbar!)" \\
\midrule
Einheit vergessen & Nur Zahl, kein J & „Energie hat IMMER die Einheit Joule (J)" \\
\midrule
Formel falsch umgestellt & $m = E \cdot g \cdot h$ & Formel-Dreieck zeigen: „Was steht oben? Was unten?" \\
\bottomrule
\end{tabularx}

% ═══════════════════════════════════════════════════════════════
% LÖSUNGEN
% ═══════════════════════════════════════════════════════════════

\newpage
\section*{Lösungen zum Arbeitsblatt}

\subsection*{Kasten 1: Wechselwirkungsgesetz}

\begin{itemize}[leftmargin=2em]
    \item Kräfte treten immer \lsg{paarweise} auf.
    \item Gleich \lsg{groß}
    \item Entgegengesetzt \lsg{gerichtet}
    \item Wirken an \lsg{verschiedenen} Körpern
\end{itemize}

\subsection*{Aufgabe 1: Kräftepaare zeichnen (6 BE)}

\begin{tabularx}{\textwidth}{@{}l|X|c@{}}
\toprule
\textbf{Situation} & \textbf{Kräftepaar} & \textbf{BE} \\
\midrule
a) Schwimmer & $F_\text{Hand→Wasser}$ (nach hinten) und $F_\text{Wasser→Hand}$ (nach vorne) & 2 \\
b) Rakete & $F_\text{Rakete→Gas}$ (nach unten) und $F_\text{Gas→Rakete}$ (nach oben) & 2 \\
c) Auto & $F_\text{Reifen→Straße}$ (nach hinten) und $F_\text{Straße→Reifen}$ (nach vorne) & 2 \\
\bottomrule
\end{tabularx}

\textit{Bewertung: Je 1 BE für korrekte Richtung, 1 BE für korrekte Beschriftung.}

\subsection*{Kasten 2: Lageenergie}

\begin{itemize}[leftmargin=2em]
    \item $E_\text{pot}$: Lageenergie in \lsg{Joule (J)}
    \item $m$: Masse in \lsg{kg}
    \item $g$: Ortsfaktor \lsg{10} m/s²
    \item $h$: Höhe in \lsg{m (Meter)}
\end{itemize}

\subsection*{Aufgabe 2: Lageenergie berechnen (6 BE)}

\begin{tabularx}{\textwidth}{@{}l|X|c@{}}
\toprule
\textbf{Teil} & \textbf{Lösung} & \textbf{BE} \\
\midrule
a) Ball & $E_\text{pot} = 2\,\text{kg} \cdot 10\,\text{m/s}^2 \cdot 1\,\text{m} = \lsg{20\,\text{J}}$ & 2 \\
b) Eimer & $E_\text{pot} = 10\,\text{kg} \cdot 10\,\text{m/s}^2 \cdot 5\,\text{m} = \lsg{500\,\text{J}}$ & 2 \\
c) Masse & $m = \frac{400\,\text{J}}{10\,\text{m/s}^2 \cdot 5\,\text{m}} = \frac{400}{50} = \lsg{8\,\text{kg}}$ & 2 \\
\bottomrule
\end{tabularx}

\subsection*{Aufgabe 3: Formel umstellen (4 BE)}

\textbf{Gegeben:} $m = 20\,\text{kg}$, $E_\text{pot} = 1000\,\text{J}$, $g = 10\,\text{m/s}^2$ \hfill (1 BE)

\textbf{Gesucht:} $h$ \hfill (0,5 BE)

\textbf{Lösung:}
\[
h = \frac{E_\text{pot}}{m \cdot g} = \frac{1000\,\text{J}}{20\,\text{kg} \cdot 10\,\text{m/s}^2} = \frac{1000}{200} = \lsg{5\,\text{m}}
\]
\hfill (2 BE für Rechnung, 0,5 BE für Antwortsatz)

\subsection*{Aufgabe 4: Staudamm (4 BE)}

\textbf{Erwartungshorizont:}

\begin{itemize}[leftmargin=2em]
    \item Wasser hat große \textbf{Masse} $m$ (1 BE)
    \item Wasser befindet sich auf einer \textbf{Höhe} $h$ über dem Tal (1 BE)
    \item Nach $E_\text{pot} = m \cdot g \cdot h$ besitzt es \textbf{Lageenergie} (1 BE)
    \item Beim Ablassen: Umwandlung in Bewegungsenergie → Turbine → Strom (1 BE)
\end{itemize}

\textit{Alternativ: Hinweis auf „je höher/mehr Wasser, desto mehr Energie" (1 BE)}

% ═══════════════════════════════════════════════════════════════
% MATERIAL-CHECKLISTE
% ═══════════════════════════════════════════════════════════════

\section*{Material-Checkliste}

\begin{itemize}[leftmargin=2em]
    \item[$\square$] Beamer + Laptop (für LP)
    \item[$\square$] AB-01b (Klassensatz, 2-seitig)
    \item[$\square$] WLAN für SuS-Geräte (optional, für LP auf Tablets)
    \item[$\square$] LP-URL / QR-Code bereit
\end{itemize}

\begin{infobox}[LP-Zugang]
\textbf{URL:} \texttt{[Server]/LP-01b-newton3-energie.html}\\
Alternativ: QR-Code auf Folie / QR-LP-01b.pdf drucken
\end{infobox}

% ═══════════════════════════════════════════════════════════════
% AUSBLICK
% ═══════════════════════════════════════════════════════════════

\section*{Ausblick}

\textbf{Nächste Stunde (01c):} Bewegungsenergie $E_\text{kin} = \frac{1}{2} m v^2$, Energieerhaltung

\textbf{Verknüpfung:} Lageenergie → Bewegungsenergie beim Fallen (Achterbahn, Pendel)

\vfill
\hrule
\vspace{0.3em}
\begin{center}
\small\textcolor{mittelgrau}{LH-01b | Stand: Januar 2026}
\end{center}

\end{document}
